% !TEX root = ../main.tex
% Бүлэг 7

\chapter{Ёс зүйн асуудал}
\label{Chapter7}
\pagecolor{Beige}

Эмнэлгийн дүрс боловсруулалтад хиймэл оюуны аргуудыг ашиглах нь оношилгооны үр ашгийг нэмэгдүүлэх боломжтой боловч өгөгдлийн хамгаалалт, шударга байдал, хариуцлагатай хэрэглээ зэрэг ёс зүйн асуудлыг давхар авч үзэх шаардлагатай байдаг. Дэлхийн эрүүл мэндийн байгууллага мөн эрүүл мэндийн хиймэл оюуныг хариуцлагатай, ил тод, хүний хяналттай ашиглах зарчмуудыг тодорхойлсон байдаг~\cite{ref45}. Энэхүү судалгаанд ашигласан өгөгдөл, загвар хөгжүүлэлт болон үр дүнгийн тайлбартай холбоотой ёс зүйн үндсэн асуудлуудыг дараах байдлаар авч үзэв.

%───────────────────────────────────────────────────────────────────────────────
\section{Өгөгдлийн зөвшөөрөл}

Судалгаанд ашигласан FracAtlas, GRAZPEDWRI-DX, YOLOBoneFracture болон BoneFractureCVProject өгөгдлийн сангууд нь олон нийтэд нээлттэй бөгөөд эрдэм шинжилгээ, судалгааны зориулалтаар ашиглах зөвшөөрөлтэй эх сурвалжууд юм. Эдгээр өгөгдлийн сангуудыг анх нийтлэх үедээ холбогдох байгууллага, судлаачид шаардлагатай зөвшөөрөл болон ёс зүйн шаардлагыг ханган боловсруулсан байдаг.

Энэхүү судалгааны хүрээнд эмнэлгийн байгууллагаас шууд өгөгдөл цуглуулаагүй бөгөөд бодит өвчтөнөөс мэдээлэл аваагүй болно. Иймээс хүний оролцоотой туршилт явуулах, өвчтөнөөс нэмэлт зөвшөөрөл авах шаардлага үүсээгүй. Судалгааны бүх өгөгдлийг зөвхөн эрдэм шинжилгээний зориулалтаар ашигласан бөгөөд эх сурвалжуудыг зохих байдлаар иш татсан.

%───────────────────────────────────────────────────────────────────────────────
\section{Нууцлал ба мэдээллийн хамгаалалт}

Эмнэлгийн дүрс нь хүний эрүүл мэндтэй холбоотой мэдрэмтгий мэдээлэл агуулж болох тул нууцлалын асуудал онцгой ач холбогдолтой байдаг. Судалгаанд ашигласан өгөгдлийн сангууд нь хувийн мэдээллээс цэвэрлэгдсэн буюу анонимжуулсан дүрсүүдээс бүрдсэн тул өвчтөнийг шууд таних боломжгүй юм~\cite{ref46}.

Судалгааны явцад өвчтөний нэр, регистрийн дугаар, холбоо барих мэдээлэл болон бусад хувийн мэдээллийг хадгалаагүй, боловсруулаагүй болно. Өгөгдлийг зөвхөн загвар сургах, үнэлэх зорилгоор ашиглаж, судалгааны хүрээнээс гадуур түгээхгүй байх зарчмыг баримталсан.

Цаашид бодит эмнэлгийн өгөгдөл ашиглах тохиолдолд өгөгдлийг анонимжуулах, байгууллагын зөвшөөрөл авах болон хувь хүний мэдээлэл хамгаалах тухай хууль, журмын холбогдох шаардлагыг хангах шаардлагатай~\cite{ref48}. Түүнчлэн сургасан загвараас сургалтын мэдээлэл задрах (model inversion) эрсдэлийг анхаарч, загвар түгээх үед зохих хамгаалалтын арга хэмжээ авах нь зүйтэй~\cite{ref49}.

%───────────────────────────────────────────────────────────────────────────────
\section{Bias ба Fairness шинжилгээ}

Хиймэл оюуны загварын үр дүн нь сургалтын өгөгдлийн бүтэц, чанараас ихээхэн хамаардаг. Судалгаанд ашигласан өгөгдлийн сангууд нь өөр өөр улс орон, эмнэлгийн байгууллагаас цуглуулсан тул хүн ам зүйн бүтэц, дүрс авах төхөөрөмж болон дүрсний чанарын хувьд ялгаатай байж болно.

Ийм нөхцөлд загвар нь сургалтын өгөгдөлд илүү түгээмэл тохиолддог хугарлын төрөл дээр сайн ажиллах боловч ховор тохиолдол дээр гүйцэтгэл буурах эрсдэлтэй. Эрүүл мэндийн салбарт ашиглагддаг алгоритмууд хүн амын бүлгүүдийн хооронд далд хазайлт үүсгэж болзошгүйг өмнөх судалгаанууд харуулсан тул энэ эрсдэлийг үнэлэх нь чухал~\cite{ref47}. Мөн ашигласан өгөгдлийн сангуудад нас, хүйс, үндэс угсаа зэрэг мэдээлэл бүрэн агуулагдаагүй тул бүлэг тус бүрийн гүйцэтгэлийг тусад нь үнэлэх боломж хязгаарлагдсан.

Бүлэг~\ref{Chapter6}-ийн алдааны шинжилгээний үр дүнгээс харахад маш жижиг хэмжээтэй, бүдэг ан цавтай хугарлууд дээр буруу ангилал харьцангуй их ажиглагдсан. Иймээс загварын үр дүнг бүх нөхцөлд ижил түвшний найдвартай гэж үзэх боломжгүй бөгөөд бодит хэрэглээнд эмчийн үнэлгээтэй хамтатган ашиглах шаардлагатай.

%───────────────────────────────────────────────────────────────────────────────
\section{Хортой хэрэглээний эрсдэл}

Эмнэлгийн зориулалтын хиймэл оюуны загваруудыг буруу ашиглах нь буруу оношилгоо, эмчилгээний хоцрогдол зэрэг сөрөг үр дагаварт хүргэж болзошгүй. Энэхүү судалгаанд боловсруулсан загвар нь туршилт, судалгааны зориулалттай бөгөөд эмнэлгийн шийдвэрийг бие даан гаргах зориулалттай биш юм.

Хэрэв загварыг клиникийн баталгаажуулалтгүйгээр шууд ашиглавал буруу ангилал эсвэл сегментчлэлийн алдаанаас шалтгаалан өвчтөний эмчилгээний шийдвэрт сөрөг нөлөө үзүүлэх эрсдэлтэй. Түүнчлэн хугарлын зэрэг үнэлгээ болон эдгэрэлтийн таамаглалын модулиуд нь туршилтын шинж чанартай тул эмчийн мэргэжлийн үнэлгээг орлох боломжгүй.

Иймээс энэхүү судалгааны үр дүнг зөвхөн судалгааны болон шийдвэр дэмжих зориулалтаар ашиглах нь зүйтэй.

%───────────────────────────────────────────────────────────────────────────────
\section{Хиймэл оюуны хэрэглээний ил тод байдал}

Энэхүү судалгаанд өгөгдөл боловсруулах, загвар хөгжүүлэх, туршилтын үр дүнг шинжлэх үйл ажиллагааг судлаач өөрөө хэрэгжүүлсэн бөгөөд хиймэл оюуны хэрэгслүүдийг зөвхөн туслах зориулалтаар ашигласан.

Бичвэрийн найруулга сайжруулах, хэл зүй болон үг үсгийн алдаа шалгах зэрэг туслах ажилд том хэлний загварт суурилсан хиймэл оюуны хэрэгслээс зөвлөгөө авсан боловч судалгааны үр дүн, туршилтын өгөгдөл, хүснэгт, зураг болон шинжилгээний үр дүнг автоматаар үүсгээгүй болно. Судалгааны бүх туршилт, загварын сургалт болон үр дүнгийн баталгаажуулалтыг судлаач өөрөө гүйцэтгэсэн.

Ил тод байдлыг хангах үүднээс судалгаанд ашигласан өгөгдлийн эх сурвалж, архитектур, туршилтын тохиргоо болон үнэлгээний аргачлалыг тайланд дэлгэрэнгүй тайлбарласан болно~\cite{ref51}. Түүнчлэн загварын шийдвэрийг тайлбарлах чадвар нь эмнэлгийн хиймэл оюуны хариуцлагатай хэрэглээний чухал бүрэлдэхүүн хэсэг тул~\cite{ref50} энэ судалгаанд Grad-CAM болон SHAP аргуудыг ашиглан загварын шийдвэрт нөлөөлж буй хүчин зүйлсийг тайлбарлахыг зорьсон.

%───────────────────────────────────────────────────────────────────────────────
\section{Бүлгийн дүгнэлт}

Энэ бүлэгт судалгаатай холбоотой ёс зүйн үндсэн асуудлуудыг авч үзэв. Судалгаанд ашигласан өгөгдөл нь олон нийтэд нээлттэй, анонимжуулсан эх сурвалжаас бүрдсэн тул хувийн мэдээллийн эрсдэл харьцангуй бага байсан. Гэсэн хэдий ч өгөгдлийн хазайлт, хүн амын төлөөлөх чадвар, хиймэл оюуны загварын алдаанаас үүдэх эрсдэлийг анхаарах шаардлагатай хэвээр байна.

Мөн боловсруулсан загвар нь эмчийн шийдвэрийг орлох бус харин дэмжих хэрэгсэл болох ёстой бөгөөд бодит эмнэлгийн орчинд ашиглахын өмнө нэмэлт баталгаажуулалт хийх шаардлагатай. Судалгааны явцад ашигласан хиймэл оюуны хэрэгслүүдийн хэрэглээг ил тод тайлагнаснаар судалгааны үнэн зөв, хариуцлагатай байдлыг хангахыг зорив.
